\chapter{Imaging Anatomy}

\label{ch:ch30_anatomy}

\section{Anatomic Orientation}
Interpretation requires knowledge of normal compartmental anatomy, planes, vascular territories, organ relationships, and expected age- or procedure-related variation. Confirm orientation and laterality before describing a lesion.

\subsection{Axial, Coronal, and Sagittal Planes}
All cross-sectional images can be understood in three orthogonal planes:
\begin{itemize}
    \item \textbf{Axial (transverse)} -- divides the body into superior and inferior halves. The standard viewing convention is from below (feet looking up), so the patient's right is on the viewer's left.
    \item \textbf{Coronal} -- divides the body into anterior and posterior halves. Useful for assessing the adrenal glands, thoracic inlet, and hip joints.
    \item \textbf{Sagittal} -- divides the body into left and right halves. Essential for spine, sternum, and midline structures such as the pituitary gland.
\end{itemize}

\subsection{Identifying Level on a Single Image}
When given a single axial image, you can determine the approximate level by identifying key landmarks:
\begin{itemize}
    \item \textbf{Supraclavicular} -- clavicle heads, lung apices, great vessels.
    \item \textbf{Aortic arch level} -- aortic arch, SVC, and trachea.
    \item \textbf{Carina} -- tracheal bifurcation, main bronchi, and nearby pulmonary arteries.
    \item \textbf{Cardiac level} -- cardiac chambers and pericardium; valve detail depends on slice position and reconstruction plane.
    \item \textbf{Liver dome} -- liver, IVC, hepatic veins.
    \item \textbf{Renal hilum} -- kidneys and renal vessels; the pancreatic body may be visible anteriorly.
    \item \textbf{Aortic bifurcation} -- common iliac vessels, sacrum.
\end{itemize}

\begin{pearlbox}[Clinical Pearl: One image, one level]
If you can identify the trachea, a great vessel, and a vertebral body, you can determine the approximate axial level within seconds. Practice this skill on every CT you read.
\end{pearlbox}

\section{Cross-Sectional Anatomy by Region}

\subsection{Brain}
On CT, grey matter is denser (lighter) than white matter. On MRI, T1-weighted images show white matter as brighter than grey matter, while T2-weighted images reverse this relationship. The ventricular system, basal ganglia, thalami, brainstem, and cerebellum should be identified on every brain study. The major vascular territories---anterior cerebral, middle cerebral, posterior cerebral, and vertebrobasilar---correspond to specific cortical and deep structures.

\subsection{Chest}
Key structures on axial CT of the chest include the trachea and main bronchi, the mediastinal great vessels (ascending aorta, descending aorta, SVC, and pulmonary arteries), the heart and pericardium, the lung parenchyma with its lobes and fissures, the pleura, the oesophagus, and the thoracic spine. On sagittal reconstructions, the aortic arch and main pulmonary artery are useful orientation landmarks.

\subsection{Abdomen}
The abdomen is divided into nine regions or four quadrants for clinical description. On axial CT, systematically identify the liver, spleen, pancreas, kidneys, adrenal glands, bowel loops, mesentery, and retroperitoneum. A portal-venous phase, often acquired around 60--70~s after contrast injection, is commonly used to assess abdominal parenchymal organs; timing and additional phases depend on the clinical question and protocol.

\subsection{Pelvis}
The pelvis contains the bladder, reproductive organs, rectum, and the pelvic floor musculature. On MRI, the zonal anatomy of the uterus (endometrium, junctional zone, myometrium) and the peripheral zone of the prostate are best assessed on T2-weighted images.

\section{Cross-Modality Correlation}
CT, MRI, ultrasound, radiography, and nuclear medicine depict different physical properties. A structure that is conspicuous on one examination may be occult on another. Correlate with the modality best suited to the clinical question.

\begin{table}[h]
\centering
\caption{Best modality for specific anatomic questions.}
\label{tab:anatomy_modality}
\begin{tabular}{p{5cm} p{4cm} p{4.5cm}}
\toprule
\textbf{Clinical Question} & \textbf{Preferred Modality} & \textbf{Why} \\
\midrule
Aortic dissection & CTA & Rapid, volumetric coverage with excellent spatial resolution \\
Bone cortex and fracture & CT or radiography & High spatial resolution for calcified structures \\
Brain parenchyma detail & MRI (T2/FLAIR/DWI) & Soft-tissue contrast resolution \\
Hip labrum & MRI; MR arthrography in selected cases & Soft-tissue assessment of the labrum and adjacent structures \\
Liver lesion characterisation & MRI with hepatobiliary contrast in selected cases & Hepatobiliary-phase features can help distinguish focal nodular hyperplasia from many adenomas, but some adenomas also retain contrast \\
Pulmonary nodules & Chest CT; low-dose for screening or follow-up when appropriate & High spatial resolution for lung parenchyma; protocol depends on the clinical question \\
\bottomrule
\end{tabular}
\end{table}

\section{Reporting}
Use precise location, side, compartment, dimensions, and relationship to critical structures. State when an organ or region is incompletely visualized. Normal anatomy should be described only to the degree necessary to support the clinical question and exclude important pathology.

\begin{pearlbox}[Clinical Pearl: Know the blood supply]
Vascular territory can help narrow the differential. A pleural-based wedge-shaped consolidation may suggest pulmonary infarction, particularly when pulmonary embolism is suspected, but it can overlap with infection and other causes. Interpret it with the full imaging pattern and clinical context.
\end{pearlbox}

\clearpage
